\documentclass{cernatlasnote}

\usepackage[T1]{fontenc}
\usepackage[utf8]{inputenc}
\usepackage{lmodern}
\usepackage{graphicx}
\usepackage{booktabs}
\usepackage{array}
\usepackage{longtable}
\usepackage{enumitem}
\usepackage{microtype}
\usepackage{xcolor}
\usepackage{needspace}

\hypersetup{
  hidelinks,
  hypertexnames=false,
  pdftitle={Public Review Draft: DUNE GIZMO--SC/DPS Interface Control Document},
  pdfauthor={Manuel Arroyave},
  pdfsubject={Interface definition for GIZMO, DUNE Slow Controls, and the Detector Protection System},
  pdfkeywords={DUNE, GIZMO, Slow Controls, Detector Protection System, OPC UA, Ethernet}
}

\setlist[itemize]{leftmargin=1.35em,itemsep=0.18em,topsep=0.28em}
\setlist[enumerate]{leftmargin=1.55em,itemsep=0.22em,topsep=0.28em}
\setlength{\parskip}{0.38em}
\setlength{\parindent}{0pt}
\renewcommand{\arraystretch}{1.12}
\setlength{\tabcolsep}{4pt}

\newcolumntype{L}[1]{>{\raggedright\arraybackslash}p{#1}}
\newcommand{\code}[1]{\nolinkurl{#1}}

\title{PUBLIC REVIEW DRAFT:\\ DUNE GIZMO -- SC/DPS\\ Interface Control Document}
\author{Manuel Arroyave}
\date{18 August 2026}
\EDMSDocNo{DUNE GIZMO--SC/DPS Interface Control Document}
\EDMSDocId{Not a controlled EDMS copy}
\draftversion{0.13-public}
\documentlabel{18 August 2026}
\DocAuthors{Manuel Arroyave}
\DocCheckedBy{TBD}
\DocApprovedBy{TBD}

\begin{document}
\maketitle

\begin{center}
\fbox{\begin{minipage}{0.91\textwidth}
\textbf{Public review copy.} This is not a controlled EDMS revision. Site
addresses, credentials, and controlled installation records are intentionally
omitted.
\end{minipage}}
\end{center}

\begin{abstract}
This ICD answers a specific integration question: how does the Ground
Impedance Monitor (GIZMO) hand its measurements, alarms, and health information
to DUNE Slow Controls and the Detector Protection System (SC/DPS)? At that
boundary, every active GIZMO looks the same: one copper Gigabit Ethernet
connection and one OPC~UA server using the \code{urn:fnal:gizmo} information
model. The Kria unit is the primary implementation and the legacy ZedBoard is
the spare. Hardware differences remain visible through OPC~UA quality and
access metadata; they do not create a second SC design. GIZMO reports the
ground-impedance alarm to SC/DPS. The response after that report is received is
defined elsewhere.
\end{abstract}

\vfill
\makereviewtable
\clearpage

\section*{Distribution List}

GIZMO system owner; DUNE Slow Controls and Detector Protection System;
installation and integration; controls cybersecurity.

\section*{History of Changes}

\begingroup
\small
\begin{longtable}{L{0.13\textwidth} L{0.16\textwidth} L{0.18\textwidth} L{0.43\textwidth}}
\toprule
\textbf{Revision} & \textbf{Date} & \textbf{Author} & \textbf{Change} \\
\midrule
0.10-public & 2026-08-14 & M. Arroyave & Dual-platform public review draft. \\
\addlinespace
0.11-public & 2026-08-17 & M. Arroyave & Replaced the platform-by-platform
design narrative with one concise interface; removed numbered GIZMO
requirements and implementation/test history; identified the ZedBoard as the
spare; clarified physical Ethernet, OPC~UA, safety-reporting, and scope
boundaries. \\
\addlinespace
0.12-public & 2026-08-17 & M. Arroyave & Renamed the document as the
GIZMO--SC/DPS ICD and revised the presentation so that context and operator
meaning accompany the normative interface requirements. \\
\addlinespace
0.13-public & 2026-08-18 & M. Arroyave & Synchronized the released OPC~UA
model baseline to 1.4.0 and clarified the boundary between GIZMO tag history
and the separate SC--Computing database service. \\
\bottomrule
\end{longtable}
\endgroup

\setcounter{tocdepth}{1}
\begingroup
\small
\tableofcontents
\endgroup
\clearpage

\section{Purpose and scope}

GIZMO watches the electrical relationship between Detector Ground and
Cavern/Building Ground. Its job at the controls boundary is straightforward:
report what it measured, report whether its own alarm logic is active, and
provide enough health information for SC/DPS to decide whether that report can
be trusted. This applies during both construction and detector operation.

The normal interface boundary is:

\begin{center}
\begin{tabular}{c}
GIZMO measurement hardware $\longrightarrow$ GIZMO controller
$\xrightarrow{\text{1000BASE-T}}$ controls network \\
controls network $\xrightarrow{\text{OPC UA}}$ SC/Ignition and SC-managed DPS
\end{tabular}
\end{center}

The ICD defines that network connection, the OPC~UA contract, the meaning and
quality of the data, the limited controls available to SC/DPS, and the tests
used to accept the interface. It includes enough grounding and rack context to
prevent the measurement from being misunderstood. It is not an electrical work
procedure, and it does not decide what happens after SC/DPS receives an alarm.
Notification, escalation, work-pause decisions, and protective actions belong
to the applicable operating and protection procedures.

In this document, ``shall'' marks behavior that is required for interface
acceptance. ``Should'' marks a recommendation. A value marked TBD is genuinely
unresolved; it is not zero, disabled, or accepted by omission.

\subsection{Applicable documents and precedence}

\begin{longtable}{L{0.32\textwidth} L{0.58\textwidth}}
\toprule
\textbf{Document} & \textbf{Why it matters at this interface} \\
\midrule
Released Far Detector grounding drawings and cable schedule & Authoritative
physical grounding topology, conductor identity, CT orientation, and
stimulus/return routing. \\
\addlinespace
Current rack ORC; DUNE-doc 35392 as the development-rack reference & Rack
configuration, power, installation, labeling, and safety authorization. The
installed rack requires its own applicable approved ORC. \\
\addlinespace
Maintained GIZMO runtime and generated
\code{schema/gizmo-opcua-contract.json} & Normative OPC~UA NodeIds, datatypes,
units, ranges, access metadata, methods, and meanings. \\
\addlinespace
DUNE-doc 1805 and 27482 & Engineering reference for the legacy ZedBoard spare;
not an alternate SC interface. \\
\addlinespace
DUNE-doc 37315 & DUNE Slow Controls architecture and integration direction. \\
\addlinespace
Current DUNE SC--Computing ICD & Database-service boundary for Slow Controls:
hosting, network access, backup, recovery, security, and support ownership. It
does not change the GIZMO OPC~UA producer contract. \\
\bottomrule
\end{longtable}

The division of authority is important. Released drawings and the applicable
ORC control the physical installation. The released generated OPC~UA contract
controls the machine interface. This ICD explains that contract for people. If
this summary or an intake spreadsheet disagrees with either controlling source,
GIZMO and SC/DPS resolve the discrepancy before the interface is used.

\section{One interface, two hardware implementations}

SC/DPS should not need a different integration simply because the electronics
inside the rack changed. It connects to a GIZMO endpoint, checks the identity
reported by that endpoint, and uses the common information model. Each active
unit runs its own OPC~UA server with its own device identity and endpoint. One
unit never proxies, configures, starts, or stops the other.

\begin{longtable}{L{0.18\textwidth} L{0.25\textwidth} L{0.47\textwidth}}
\toprule
\textbf{Role} & \textbf{Hardware} & \textbf{SC/DPS treatment} \\
\midrule
Primary & AMD/Xilinx K26 SOM on a KR260 carrier & Full reference implementation
of the common interface. \\
\addlinespace
Spare & Legacy ZedBoard and preserved controller & Same namespace, NodeIds,
datatypes, units, and meanings. Unsupported data remain visible with a
specific non-good StatusCode; unsupported controls return
\code{BadNotSupported}. A narrower hardware limit may reject a write without
changing the common engineering metadata. \\
\bottomrule
\end{longtable}

The hardware name remains useful for maintenance and may be shown on the
display. It must not change the meaning of a tag. Dashboards, alarms, history,
and operator workflow use the common model, so changing to the spare does not
require a second SC/DPS tag design.

\section{Physical interface}

\subsection{SC/DPS network connection}

For SC/DPS, the physical handoff is one ordinary copper Ethernet link per
active GIZMO. The details below make that statement testable and prevent a
maintenance port from quietly becoming a second controls interface.

\Needspace{0.34\textheight}
\begin{longtable}{L{0.25\textwidth} L{0.65\textwidth}}
\toprule
\textbf{Item} & \textbf{Interface} \\
\midrule
Link & 1000BASE-T Gigabit Ethernet in accordance with IEEE 802.3. The accepted
operating state is 1~Gb/s, full duplex. \\
Medium & Four-pair balanced copper cable, Category 5e or better, installed to
the site cabling standard. The complete channel, including patch cables, shall
not exceed 100~m. \\
Connector & 8P8C modular connector, commonly called RJ-45. \\
Negotiation & Auto-negotiation may be used; link speed and duplex shall be
verified during commissioning. \\
Network & Site-assigned hostname, address, switch port, VLAN, route, time
source, and firewall policy on the restricted controls network. These are
installation records, not fixed in this public ICD. \\
Labeling & Both cable ends carry the device and switch-port identifiers.
Label the rack with its equipment identifier and the applicable ORC number. \\
\bottomrule
\end{longtable}

A second controller port, a maintenance laptop connection, USB, UART, JTAG,
and local web pages are maintenance interfaces. They are not additional SC
data or command paths.

\subsection{Rack grounding and safety context}

The network interface cannot be interpreted safely without knowing what the
instrument is watching. Rack-door and endpoint labels therefore identify
Detector Ground, Cavern/Building Ground, the monitored safety-ground conductor
through the current transformer (CT), and the approved stimulus and return
points. If the released installation adopts the site color convention, Cavern
Ground is green and Detector Ground is green with a yellow stripe. Color is a
visual aid, never the source of truth; the released drawing and durable text
labels remain authoritative.

\begin{center}
\fbox{\begin{minipage}{0.89\textwidth}
\underline{\textbf{The detector is not simply left floating.}} An approved
safety-ground path connects the two references through the saturable-inductor
assembly. GIZMO observes that arrangement; it is not the safety ground.

\underline{\textbf{Never lift the safety ground}} to clear a GIZMO alarm or
restore a reading.
\end{minipage}}
\end{center}

The released electrical design and ORC own the CT, stimulus, power, chassis
bond, beacon/audio, and saturable-inductor connections. They are context for
SC/DPS, not controls exposed by this ICD. SC/DPS must not alter or drive them
through an undocumented path.

The rack door also carries a short, revision-controlled reference identifying
the rack and device, the applicable ORC, both ground references, the meaning of
the local alarm indications, and the safety statement above. That reference
may point the reader to the controlled response procedure; it does not invent
or summarize SC/DPS escalation logic.

\section{Protocol}

Once the Ethernet link is up, SC/DPS talks to GIZMO through OPC~UA. Each active
GIZMO shall expose one OPC~UA server as its supported public machine interface;
the client does not need to know how the controller moves data internally.

\begin{longtable}{L{0.25\textwidth} L{0.65\textwidth}}
\toprule
\textbf{Item} & \textbf{Interface} \\
\midrule
Protocol & OPC~UA Binary over UA TCP. \\
Endpoint & \code{opc.tcp://<assigned-host>:<assigned-port>}; TCP 4840 is the
default. Each device has a distinct endpoint and application/device identity. \\
Namespace & \code{urn:fnal:gizmo}. Clients resolve the namespace URI at
connection time and shall not hard-code its numeric namespace index. \\
Root object & \code{Objects/GIZMo}. \\
Contract & Released generated GIZMO OPC~UA JSON schema, currently model 1.4.0.
Stable string NodeIds, datatypes, ranks, units, ranges, access metadata, method
signatures, quality, and timestamp meanings form the contract. \\
Client behavior & SC/DPS uses persistent OPC~UA subscriptions. Repeatedly opening
sessions or polling private device services is not supported. \\
Security & The endpoint resides on the restricted controls network. The
approved site certificate, trust-list, role, and firewall policy shall be
recorded at installation. An insecure policy requires an explicit approved
exception and is not encryption. \\
\bottomrule
\end{longtable}

Private device protocols, compatibility ports, SSH, HTTP dashboards, service
managers, and raw register access may still be useful to a maintainer. They are
not part of this interface, and an SC/DPS client must not depend on them.

\section{Data and controls}

\subsection{Minimum SC/DPS view}

The generated contract contains the complete address space. An operator does
not need to see every node at once. The normal SC/DPS view should instead make
three things immediately clear: which GIZMO is reporting, what it currently
sees, and whether that report is trustworthy. The groups below are the minimum
needed to answer those questions; exact NodeIds and metadata remain in the
generated schema.

\begin{longtable}{L{0.20\textwidth} L{0.31\textwidth} L{0.39\textwidth}}
\toprule
\textbf{Group} & \textbf{Essential values} & \textbf{Purpose} \\
\midrule
Identity & DeviceId, product/platform, model version, runtime version, boot ID
& Answers ``Which instrument and software produced this report?'' \\
Measurement & Sequence, source time, resistance, resistance range, active
threshold, stimulus frequency, quality, diagnostic & Shows the measured state
and whether samples are continuing to arrive. \\
Alarm & Active, Latched, Reason, LatchTime & Shows the device's authoritative
ground-impedance/ground-short decision and persistent latch. \\
Configuration & ThresholdOhm, AveragesPerCalculation, LastCommandResult &
Provides bounded configuration and command readback. \\
Calibration & State and active calibration identity/time & Distinguishes valid,
unvalidated, and maintenance states. \\
Time & Current UTC and synchronization status & Establishes trustworthy source
timestamps. \\
Health & Overall, measurement, network, services, firmware, and last update &
Separates device/interface faults from a physical ground alarm. \\
\bottomrule
\end{longtable}

Measurement and alarm values normally update once per second. Platform health
changes more slowly and may update every 10 or 30 seconds. Each sample carries
the time and sequence assigned at its source. SC/DPS preserves those values and
does not fill a gap with invented intermediate samples.

\subsection{Quality and value semantics}

The number alone is never the whole report. Every value arrives with an OPC~UA
datatype, a StatusCode, and a source timestamp. Together they tell SC/DPS both
what GIZMO measured and how much confidence to place in it. Unknown or
unavailable numeric data must not masquerade as zero. Before the first sample,
the value reports a waiting or bad state. If the source is later lost, the last
value may remain visible, but its status becomes uncertain or bad so that it
cannot look current.

Resistance illustrates why those fields must be read together.
\code{ResistanceOhm} is numeric only when
\code{ResistanceRange=InRange}. Above the validated presentation range, GIZMO
reports \code{ResistanceRange=OutOfRange}; the display calls this
\code{HIGH Z}, and the resistance itself is non-numeric. A good-quality
HIGH~Z report is a valid range state. It is neither a sensor fault nor an
invented 500- or 999-ohm measurement.

\code{Alarm.Active} comes from the device logic that drives the physical alarm
output. That Boolean is the alarm decision SC/DPS uses. Recreating the decision
with a separate resistance-versus-threshold comparison could disagree with the
instrument and is therefore not allowed. The report says that the configured
alarm condition exists; it cannot locate or diagnose the conductive path.

\subsection{Controls boundary}

Most information crosses this interface from GIZMO to SC/DPS. The few controls
that travel in the other direction are deliberately bounded and must always
return an observable result.

\begin{longtable}{L{0.31\textwidth} L{0.21\textwidth} L{0.38\textwidth}}
\toprule
\textbf{Control} & \textbf{Class} & \textbf{Rule} \\
\midrule
ThresholdOhm & Bounded read/write & Authorized operator setting;
verify the live measurement readback after writing. The spare may enforce a
narrower accepted range and return \code{BadOutOfRange}. \\
AveragesPerCalculation & Bounded read/write when supported &
Verify readback and a new measurement sequence. An unsupported spare write
returns \code{BadNotSupported}. \\
ClearLatch & Operator method when supported & Clears only the
persistent latch; it shall not suppress an active physical alarm. \\
Calibration, ADC capture, normalization, and manual time operations &
Maintenance-only & Disabled in the normal SC view until authorization,
readback, failure cleanup, and commissioning are approved. \\
\bottomrule
\end{longtable}

Every write or method call is typed, bounded, role-authorized, serialized,
logged, and followed by a result or readback. A successful method return means
that GIZMO accepted the request; the resulting measurement sequence or state
change is what demonstrates completion. Shell commands, service commands,
reboots, file paths, FPGA/register writes, relay manipulation, and raw
calibration-resistor selection are maintenance actions, not SC/DPS controls.

\section{How SC/DPS uses the report}

\subsection{Display, alarms, and history}

An operator first needs to know whether there is a ground alarm, but a clear
alarm is meaningful only when the underlying report is current and valid. The
main display therefore shows device identity, endpoint health, data age,
StatusCode/quality, resistance or HIGH~Z, threshold, alarm/latch/reason, and
calibration state together.

SC/DPS presents the following as distinct conditions rather than collapsing
them into a single generic GIZMO alarm:

\begin{itemize}
  \item the physical ground-impedance alarm reported by
  \code{GIZMo.Alarm.Active};
  \item communication loss or a stale timestamp/sequence;
  \item non-good measurement quality or invalid calibration;
  \item unexpected device/model identity; and
  \item a failed or rejected command.
\end{itemize}

The stale-data limit is selected and approved during commissioning. Crossing
that limit raises a monitoring-coverage alarm. It does not turn the last value
into a good report, and it must not be mistaken for a clear ground state. The
operating procedure---not this ICD---decides what activity may continue while
coverage is unavailable.

Ignition subscribes directly to the GIZMO OPC~UA endpoint. Its history retains
the device identity, source timestamp, StatusCode, range state, alarm
transitions, and command audit throughout construction and operation. When the
ground-impedance report reaches the SC-managed DPS integration, it still carries
the quality and source timestamp supplied by GIZMO.

The database behind Ignition is not a GIZMO device-interface requirement. The
selected long-term direction is PostgreSQL, while hosting, network access,
retention, backup, recovery, security, and service ownership follow the
SC--Computing ICD and site operations. An interim test database does not define
the production service.

That is the end of this interface: SC/DPS has received, quality-checked,
displayed, alarmed, and recorded the report. Notification, escalation,
work-control, and automatic or manual protective response begin on the other
side of that boundary and are defined elsewhere.

\section{Roles and responsibilities}

Responsibility follows the report: GIZMO owns it at the source; SC/DPS owns its
secure delivery, presentation, freshness supervision, and recording. This
allocation stops at the reporting boundary. The downstream response is defined
elsewhere.

\textbf{The GIZMO system owner} owns the instrument, common OPC~UA contract,
device software, value meanings, calibration authority, and safe local alarm
behavior. The owner supplies the device identity, data, quality, and source
timestamps.

\textbf{SC/DPS} owns the secure OPC~UA integration, operator display,
supervisory and stale-data alarms, Ignition history, command audit, and delivery
of the GIZMO report to DPS. Computing service responsibilities remain allocated
by the separate SC--Computing ICD.

\textbf{Together, GIZMO and SC/DPS} commission the endpoint, map the signals,
accept the interface, and control later changes.

\section{Interface acceptance}

Acceptance follows a report from the device to the operator display and DPS
handoff, then checks the few controls in the reverse direction. Before the
interface is released, GIZMO and SC/DPS must jointly complete these checks:

\begin{enumerate}
  \item verify the rack ID and applicable ORC label, cable labels, assigned
  switch port, 1~Gb/s full-duplex link, VLAN/addressing, route, and time source;
  \item connect with a generic OPC~UA client, resolve
  \code{urn:fnal:gizmo}, and verify device/application identity;
  \item run the versioned schema conformance test for required NodeIds,
  datatypes, units, access, StatusCodes, and source timestamps;
  \item verify once-per-second measurement sequence/time progression and the
  commissioned stale-data alarm;
  \item with an approved test input, verify finite resistance, HIGH~Z, the
  device alarm state, latch, reason, bad or unavailable data, disconnect,
  restart, and
  recovery without fabricated values;
  \item verify that Ignition and the SC-managed DPS mapping use
  \code{Alarm.Active}, preserve source quality and time in history, and record
  command results/readbacks;
  \item reject anonymous or unauthorized writes and every out-of-range or
  unsupported operation explicitly; and
  \item demonstrate that loss of OPC~UA or Ignition history does not stop local
  measurement or the physical alarm.
\end{enumerate}

The spare passes the same interface tests before it is declared available. A
successful interface test demonstrates that data and controls cross the
boundary correctly. It does not establish electrical authorization, metrology,
calibration validity, or approval of any response downstream of SC/DPS.

\section{Items requiring controlled approval}

The interface can be described publicly without publishing site configuration.
The decisions below are therefore completed in controlled records before
deployment; none may be filled in by assumption or inherited silently from a
prototype.

\Needspace{0.28\textheight}
\begingroup
\small
\begin{longtable}{L{0.60\textwidth} L{0.30\textwidth}}
\toprule
\textbf{Decision} & \textbf{Owner} \\
\midrule
Assign the EDMS number, reviewers, approvers, and controlled distribution. &
GIZMO + SC/DPS \\
Assign the underground rack ORC and place its number on the rack label. &
Installation + electrical safety \\
Approve production endpoint, switch/VLAN, addressing, certificates, roles,
and firewall rules. & SC + cybersecurity \\
Approve site threshold, stale-data limit, alarm priority, and SC/DPS signal
mapping. & GIZMO + SC/DPS \\
Approve the rack-door reference content and revision control. & GIZMO + SC/DPS +
electrical safety \\
Complete common-interface and physical-display acceptance of the ZedBoard
spare. & GIZMO + SC/DPS \\
\bottomrule
\end{longtable}
\endgroup

\end{document}
